% Generated draft: requires mathematical and source review.
\documentclass[11pt]{article}
\usepackage[margin=25mm]{geometry}
\usepackage{fontspec}
\setmainfont{Latin Modern Roman}
\usepackage{amsmath,amssymb,mathtools,unicode-math}
\setmathfont{Latin Modern Math}
\usepackage{enumitem,needspace,xurl,hyperref}
\setlength{\parindent}{0pt}
\setlength{\parskip}{5pt}
\setlength{\emergencystretch}{3em}
\newcommand{\N}{\mathbb{N}}
\newcommand{\Z}{\mathbb{Z}}
\newcommand{\Q}{\mathbb{Q}}
\newcommand{\R}{\mathbb{R}}
\newcommand{\C}{\mathbb{C}}
\newcommand{\F}{\mathbb{F}}
\newcommand{\A}{\mathbb{A}}
\newcommand{\PP}{\mathbb{P}}
\newcommand{\E}{\mathbb{E}}
\newcommand{\eps}{\varepsilon}
\newcommand{\dd}{\,\mathrm{d}}
\begin{document}
\section*{UnsolvedMath ID: 30002822}
\textbf{Source problem code:} catalogue entry
\subsection{Definitions and mathematical statement}
For the usual simple wheel graph \(W_n\), with \(n\geq 4\), the vertices consist of a cycle on \(n-1\) vertices and one additional vertex joined to every cycle vertex. Thus \(|V(W_n)|=n\) and \(|E(W_n)|=2n-2\). This describes the standard simple-wheel convention; the supplied record does not explicitly specify its admissible values of \(n\) or graph conventions for the proposed generalization.

A self-stress ordinarily refers to a geometric framework, not just an abstract graph. For example, for a framework \(p:V(G)\to\R^d\), a real equilibrium self-stress is an assignment \(\omega:E(G)\to\R\) satisfying
\[
\sum_{j:\,ij\in E(G)}\omega_{ij}(p_j-p_i)=0
\qquad\text{for every }i\in V(G).
\]
This is a conventional definition supplied to identify the missing geometric data. The record does not specify that this is the intended convention, the dimension \(d\), the placement \(p\), or its required properties. In particular, the word ``planar'' does not determine whether a specific straight-line realization is intended.

For scalar edge assignments \(\omega=(\omega_{ij})\) and \(f=(f_{ij})\) over a common field, write their edge pairing as
\[
\langle\omega,f\rangle_{E(G)}
:=\sum_{ij\in E(G)}\omega_{ij}f_{ij},
\]
where each undirected edge is counted once. The record supplies no scalar field or definition of \(f_{ij}\).

The supplied wheel assertion is that the intended self-stress and edge function satisfy
\[
\langle\omega,f\rangle_{E(W_n)}=1.
\]
The requested problem is to determine which other graphs \(G\) with
\[
|V(G)|=n,\qquad |E(G)|=2n-2
\]
admit the corresponding self-stress identity, and to identify a general phenomenon, possibly homological or cohomological, that explains it.

This is an unresolved formulation rather than a fully specified classification problem: ``corresponding'' depends on the omitted construction of \(\omega\) and \(f\), their normalization, the allowed framework realizations, and the quantifiers governing those realizations. No particular chain complex, homology theory, or cohomology theory is specified. Consequently, the supplied evidence does not determine a precise predicate on graphs or a specific homological statement. The wheel assertion and the intended generalization are retained as supplied, without asserting a theorem under additional hypotheses.
\subsection{Short English statement}
The supplied problem asks which graphs with \(n\) vertices and \(2n-2\) edges share a stated wheel self-stress identity, and whether homology or cohomology explains it. The framework, edge function, normalization, and realization quantifiers are missing, so the intended classification cannot yet be stated precisely.
\subsection{Sources}
\begin{itemize}
\item Local registry: \texttt{data/unsolvedmath/problems.json}, supplied record with permanent UnsolvedMath ID \(30002822\), title ``Self-Stress Identities for Sparse Graphs,'' and exact source problem code \texttt{catalogue entry}. The record belongs to the Oberwolfach Reports Open Problems collection and gives proposed year 2015; no individual proposer is supplied.
\item Supplied report citation: \emph{Discrete Differential Geometry} (2015), \emph{Oberwolfach Reports}, DOI 10.4171/owr/2015/13, pp.~661--729. Source URL: \url{https://doi.org/10.4171/owr/2015/13}. The supplied background describes Oberwolfach Reports as documenting mathematical workshops and their open-problem discussions. No problem-specific page or full report passage is supplied.
\item Supplied literature source: ``OWR Report 13/2015,'' \url{https://doi.org/10.4171/owr/2015/13}. The literature assessment dated \texttt{2026-08-21} labels the problem open and states that the wheel self-stress identity is known, while the requested classification and its homological or cohomological meaning remain unresolved. This is supplied dated curation derived from the OWR material, not an independently verified literature assessment; its status and cited claims require independent verification.
\item The registry labels the statement \texttt{corrected\_verified} and reports a manual semantic, typographic, and LaTeX review claiming a self-contained, meaning-preserving cleanup. These are supplied editorial assertions, not independent verification of the missing mathematical definitions or of mathematical correctness.
\end{itemize}
\end{document}
